%% =====================================================================
%%  Corrigés — Unit 15 « Let's go out! »
%% =====================================================================
\corrigesUnite{15}{Let's go out!}

\begin{corrige}{1}
\begin{reponses}[4]
\item \rep{a musical}
\item \rep{a ticket}
\item \rep{a party}
\item \rep{go to the movies}
\item \rep{busy}
\item \rep{free}
\item \rep{tonight}
\item \rep{go out}
\end{reponses}
\retiens{\en{free} a deux sens : « libre » (\en{Are you free?}) et « gratuit » (\en{free tickets}).}
\end{corrige}

\begin{corrige}{2}
\begin{reponses}[1]
\item \rep{Call him tonight!}
\item \rep{I can see her.}
\item \rep{Come with us!}
\item \rep{I don't like it.}
\item \rep{Can you help them?}
\item \rep{Mrs.\ Lopez invites her.}
\end{reponses}
\pourquoi{Sam → \en{him} ; Emma, Léa → \en{her} ; Jack et moi → \en{us} ; le film → \en{it} ; Léa et
Sam → \en{them}. Le pronom se place après le verbe ou la préposition.}
\end{corrige}

\begin{corrige}{3}
\begin{reponses}[6][french]
\item \rep{A}
\item \rep{R}
\item \rep{A}
\item \rep{R}
\item \rep{A}
\item \rep{R}
\end{reponses}
\retiens{un refus poli contient presque toujours une excuse (\en{I'm sorry}) et une raison (\en{I'm
busy}) ou une ouverture (\en{Maybe another time}).}
\end{corrige}

\begin{corrige}{4}
\begin{reponses}[1]
\item \rep{Let's have dinner!}
\item \rep{Would you like to come to my party?}
\item \rep{Let's meet at seven.}
\item \rep{Would you like to see a musical?}
\item \rep{Let's go out tonight!}
\end{reponses}
\piege{n'oubliez pas \en{to} après \en{Would you like}, et ne le mettez jamais après \en{Let's}.}
\end{corrige}

\begin{corrige}{5}
\begin{reponses}[3][french]
\item \rep{F}
\item \rep{V}
\item \rep{F}
\item \rep{V}
\item \rep{V}
\item \rep{F}
\end{reponses}
\pourquoi{la fête a lieu mercredi (\en{Wednesday}), la veille de Thanksgiving, dans la cuisine du
studio. Chacun apporte un plat de son pays. La surprise est pour Jack.}
\end{corrige}

\begin{corrige}{6}
\begin{reponses}[6][french]
\item \rep{b}
\item \rep{c}
\item \rep{e}
\item \rep{a}
\item \rep{d}
\item \rep{f}
\end{reponses}
\end{corrige}

\begin{corrige}{7}
\begin{reponses}[3][french]
\item \rep{accepte}
\item \rep{refuse}
\item \rep{accepte}
\item \rep{refuse}
\item \rep{accepte}
\item \rep{refuse}
\end{reponses}
\pourquoi{au 4, Mr.\ Carter refuse sans dire \en{no} : \en{That's very kind, but…} (« C'est très gentil,
mais… ») annonce toujours un refus. Texte de la piste page~\pageref{trans-15.10}.}
\end{corrige}

\begin{corrige}{8}
\begin{reponses}[5]
\item \rep{Saturday}
\item \rep{museum}
\item \rep{Chinatown}
\item \rep{London}
\item \rep{call}
\end{reponses}
\pourquoi{texte de la piste page~\pageref{trans-15.11}. \en{Call me back!} = « Rappelle-moi ! ».}
\end{corrige}

\begin{corrige}{9}
\begin{reponses}[4]
\item \rep{free}
\item \rep{like}
\item \rep{love}
\item \rep{Where}
\item \rep{Let's}
\item \rep{meet}
\item \rep{me}
\item \rep{it}
\end{reponses}
\retiens{\en{text it to you} = « te l'envoyer par texto » : \en{text} est aussi un verbe.}
\end{corrige}

\begin{corrige}{10}
\begin{reponses}[1]
\item \rep{Are you free on Saturday? — Yes, I am.}
\item \rep{Would you like to come to my party? — I'd love to!}
\item \rep{Let's go to the movies tonight!}
\item \rep{I'm sorry, I can't. I'm busy.}
\item \rep{Let's meet at eight, in front of the theater.}
\item \rep{Call me tomorrow!}
\item \rep{I like Sam. I invite him.}
\item \rep{Come with us!}
\end{reponses}
\piege{phrase 7 : « je l'invite » → \en{I invite \textbf{him}} : le pronom se place après le verbe.}
\end{corrige}

\begin{corrige}{11}
Réponse modèle :
\begin{anglais}
\begin{tcolorbox}[otBase,colback=otGreenTint,colframe=otGreenTint,fontupper=\fontsize{10.4}{14}\selectfont]
Hi Tom! Are you free on Saturday? Would you like to go to the movies with me? There's a great new
movie at the Rex. Let's meet at seven, in front of the theater. Then let's have a pizza! Text me!
\end{tcolorbox}
\end{anglais}
\end{corrige}

\begin{corrige}{12}
Réponses modèles :
\begin{reponses}[1]
\item \rep{Yes, I am.}
\item \rep{I'd love to!} \textit{ou} \rep{That sounds great!}
\item \rep{Sure! Where?}
\item \rep{I'm sorry, I can't. Maybe another time.}
\item \rep{Sure! Call me tomorrow.}
\end{reponses}
\end{corrige}

\begin{corrigetest}
\begin{reponses}[5][french]
\item \rep{b}
\item \rep{a}
\item \rep{c}
\item \rep{b}
\item \rep{c}
\item \rep{a}
\item \rep{b}
\item \rep{a}
\item \rep{c}
\item \rep{b}
\end{reponses}
\pourquoi{question 1 : \en{Do you like…?} demande un goût, pas une invitation. Question 8 : à une
question fermée, on répond \en{yes} ou \en{no} ; la voix monte.}
\end{corrigetest}
